\documentclass{article}
\usepackage[utf8]{inputenc}
\usepackage{amsmath}
\usepackage{geometry}
\geometry{margin=2cm}
\begin{document}

\section*{Análise da Questão 177 — ENEM 2024 — Matemática}

\subsection*{Enunciado}
Uma empresa produz mochilas escolares sob encomenda. O custo total de produção é composto por um custo fixo, que não depende do número de mochilas, e um custo variável, proporcional à quantidade produzida. O custo total cresce de forma linear. A tabela abaixo apresenta o custo total para três quantidades distintas de mochilas produzidas.

\begin{center}
\begin{tabular}{|c|c|c|c|}
  \hline
  \rowcolor{pink}
  \textbf{Quantidade de mochilas} & 30 & 50 & 100 \\
  \hline
  \textbf{Custo total (R\$)} & 1050,00 & 1650,00 & 3150,00 \\
  \hline
\end{tabular}
\end{center}

O custo total para produzir 80 mochilas será solicitado.

\section*{Solução Técnica}

Considere que o custo total $C$ em reais depende da quantidade de mochilas produzidas $x$ segundo uma função linear:

\[
C = ax + b,
\]

onde $a$ representa o custo variável por mochila e $b$ o custo fixo.

\textbf{Passos:}
\begin{enumerate}
  \item Dados os pontos $(x, C)$ da tabela: (30, 1050), (50, 1650) e (100, 3150).
  \item Calcule o coeficiente angular $a$ com os pontos $(30,1050)$ e $(50,1650)$:
  \[
  a = \frac{1650 - 1050}{50 - 30} = \frac{600}{20} = 30.
  \]
  \item Encontre o custo fixo $b$ usando o ponto $(30, 1050)$:
  \[
  1050 = 30 \times 30 + b \implies 1050 = 900 + b \implies b = 150.
  \]
  \item Portanto, a função do custo total é:
  \[
  C = 30x + 150.
  \]
  \item Calcule o custo para $x=80$ mochilas:
  \[
  C = 30 \times 80 + 150 = 2400 + 150 = 2550.
  \]
\end{enumerate}

\textbf{Resposta:} O custo total para a produção de 80 mochilas será \textbf{R\$ 2\,550,00}.

\section*{Fórmulas Utilizadas}

\begin{itemize}
  \item Função linear do custo total:
  \[
  C = ax + b,
  \]
  onde $C$ é o custo total, $x$ a quantidade produzida, $a$ o custo variável por unidade, e $b$ o custo fixo.

  \item Coeficiente angular (taxa de variação):
  \[
  a = \frac{C_2 - C_1}{x_2 - x_1}.
  \]
\end{itemize}

\section*{Evidência Visual Utilizada}

A tabela fornecida apresenta os seguintes dados importantes para a formulação da função linear:

\begin{center}
\begin{tabular}{|c|c|c|c|}
  \hline
  \rowcolor{pink}
  \textbf{Quantidade de mochilas} & 30 & 50 & 100 \\
  \hline
  \textbf{Custo total (R\$)} & 1050,00 & 1650,00 & 3150,00 \\
  \hline
\end{tabular}
\end{center}

\section*{Explicação Didática}

\subsection*{Resumo}

Esta questão trabalha o conceito de função linear para modelar o custo total de produção, que é a soma de um custo fixo e de um custo variável proporcional à quantidade produzida. A partir dos dados da tabela, construímos a função e depois a utilizamos para calcular um custo intermediário.

\subsection*{Passo a Passo}

\begin{enumerate}
  \item Observe os pontos dados na tabela: (30, 1050), (50, 1650), (100, 3150).
  \item Identifique que $C = ax + b$, uma função linear, descreve o custo total.
  \item Calcule o coeficiente angular:
  \[
  a = \frac{1650 - 1050}{50 - 30} = 30,
  \]
  ou seja, cada mochila custa R\$30 adicionais.
  \item Determine o custo fixo $b$ usando um ponto da tabela:
  \[
  1050 = 30 \times 30 + b \Rightarrow b = 150.
  \]
  \item A função é então:
  \[
  C = 30x + 150.
  \]
  \item Calcule para 80 mochilas:
  \[
  C = 30 \times 80 + 150 = 2\,550.
  \]
  \item Portanto, o custo para 80 mochilas é R\$2\,550,00.
\end{enumerate}

\subsection*{Erros Comuns}

\begin{itemize}
  \item Ignorar o custo fixo e multiplicar apenas o custo variável.
  \item Calcular errado o coeficiente angular.
  \item Confundir o custo fixo com a taxa de variação.
  \item Interpretação incorreta da tabela.
  \item Usar média simples sem modelar com função.
\end{itemize}

\subsection*{Dicas para Ensino}

\begin{itemize}
  \item Demonstre a diferença entre custos fixos e variáveis.
  \item Ensine a calcular taxa de variação e formação da função linear.
  \item Explore tabelas e gráficos para visualização.
  \item Pratique exercícios de interpolação e extrapolação.
\end{itemize}

\section*{Distratores}

\begin{itemize}
  \item \textbf{R\$2\,400,00} - Multiplicar apenas custo variável $30\times 80$ e esquecer o custo fixo.
  \item \textbf{R\$2\,520,00} - Erro no cálculo do coeficiente angular ou cálculo do custo fixo.
  \item \textbf{R\$2\,700,00} - Confusão na interpretação dos pontos e proporções.
  \item \textbf{R\$2\,800,00} - Erro de soma ou interpretação incorreta dos dados da tabela.
\end{itemize}

\section*{Perfil do Item}

\textbf{Habilidade Primária:} Modelagem e interpretação de funções lineares.

\textbf{Habilidades Secundárias:}
\begin{itemize}
  \item Cálculo do coeficiente angular.
  \item Interpretação de tabelas.
  \item Aplicação prática da matemática.
\end{itemize}

\textbf{Demanda Cognitiva:} Média.

\textbf{Dificuldade:} Média.

\textbf{Requisitos: Visualização, modelagem e cálculo.}

\textbf{Notas:}
\begin{itemize}
  \item Entender custo fixo e variável.
  \item Interpretar dados tabelados.
  \item Calcular coeficiente angular e aplicar substituição.
\end{itemize}

\end{document}